% !TeX program = xelatex
\documentclass[UTF8,11pt,oneside,openany,fontset=none]{ctexbook}
\usepackage[paperwidth=180mm,paperheight=255mm,inner=23mm,outer=23mm,top=24mm,bottom=23mm,headheight=23pt,headsep=8mm,footskip=11mm,marginparwidth=15mm,marginparsep=3mm]{geometry}
\usepackage{fontspec}
\setmainfont{texgyrepagella-regular.otf}[BoldFont=texgyrepagella-bold.otf,ItalicFont=texgyrepagella-italic.otf,BoldItalicFont=texgyrepagella-bolditalic.otf]
\setsansfont{texgyreheros-regular.otf}[BoldFont=texgyreheros-bold.otf,ItalicFont=texgyreheros-italic.otf]
\newfontfamily\cyrillicfont{Times New Roman}
\setmonofont{lmmono10-regular.otf}[Scale=MatchLowercase]
\setCJKmainfont{Songti SC}[BoldFont={Songti SC Bold},ItalicFont={Kaiti SC}]
\setCJKsansfont{Heiti SC}
\setCJKmonofont{Kaiti SC}
\usepackage{amsmath,amssymb,mathrsfs,mathtools}
\usepackage{graphicx,xcolor,tikz,tikz-cd,booktabs,longtable,array,enumitem,etoolbox,needspace}
\usetikzlibrary{arrows.meta,calc,decorations.markings,patterns}
\usepackage[bottom,hang]{footmisc}
\usepackage{fancyhdr}
\usepackage{microtype}
\definecolor{Ink}{HTML}{18323D}
\definecolor{Teal}{HTML}{266A70}
\definecolor{Gold}{HTML}{C5A363}
\definecolor{Muted}{HTML}{63777E}
\definecolor{Paper}{HTML}{F8F4EC}
\usepackage[unicode,bookmarksnumbered=true,bookmarksopen=true,bookmarksopenlevel=1,colorlinks=true,linkcolor=Teal,citecolor=Teal,urlcolor=Teal,pdfborder={0 0 0},pdfencoding=auto]{hyperref}
\usepackage{bookmark}
\hypersetup{pdftitle={函子代数几何导论·第一卷：仿射代数几何},pdfauthor={A. Grothendieck（讲授）；F. Gaeta（笔记）},pdfsubject={完整中文译本；中英法术语对照；XeLaTeX}}
\ctexset{chapter={name={第,章},number=\arabic{chapter},format=\raggedright\sffamily\bfseries\color{Ink}\Huge,beforeskip=18pt,afterskip=28pt},section={format=\sffamily\bfseries\color{Ink}\Large,beforeskip=24pt,afterskip=12pt},subsection={format=\sffamily\bfseries\color{Teal}\large,beforeskip=20pt,afterskip=10pt}}
\setlength{\parindent}{2em}
\setlength{\parskip}{0.22em plus 0.12em minus 0.06em}
\linespread{1.23}
\setlength{\emergencystretch}{2em}
\allowdisplaybreaks[2]
\setlist{nosep,leftmargin=2em,topsep=0.5em,itemsep=0.25em}
\renewcommand{\arraystretch}{1.22}
\setcounter{tocdepth}{1}
\setcounter{secnumdepth}{0}
\pagestyle{fancy}
\fancyhf{}
\fancyhead[L]{\color{Muted}\footnotesize\sffamily 函子代数几何导论\enspace /\enspace 仿射代数几何}
\fancyhead[R]{\color{Muted}\footnotesize\nouppercase{\leftmark}}
\fancyfoot[C]{\color{Teal}\small\thepage}
\renewcommand{\headrulewidth}{0.3pt}
\renewcommand{\headrule}{\hbox to\headwidth{\color{Gold}\leaders\hrule height\headrulewidth\hfill}}
\fancypagestyle{plain}{\fancyhf{}\fancyfoot[C]{\color{Teal}\small\thepage}\renewcommand{\headrulewidth}{0pt}}
\makeatletter
\renewcommand{\chaptermark}[1]{\if@mainmatter\markboth{第\thechapter 章\quad #1}{}\else\markboth{#1}{}\fi}
\makeatother
\DeclareMathOperator{\Hom}{Hom}
\DeclareMathOperator{\Spec}{Spec}
\DeclareMathOperator{\Aff}{Aff}
\DeclareMathOperator{\Sets}{Sets}
\DeclareMathOperator{\id}{id}
\newcommand{\op}{\mathrm{op}}
\DeclareMathOperator*{\colim}{colim}
\DeclareMathOperator{\coker}{coker}
\DeclareMathOperator{\Ann}{Ann}
\DeclareMathOperator{\Frac}{Frac}
\DeclareMathOperator{\Supp}{Supp}
\DeclareMathOperator{\Max}{Max}
\DeclareMathOperator{\End}{End}
\DeclareMathOperator{\Aut}{Aut}
\DeclareMathOperator{\Ker}{Ker}
\let\Im\relax
\DeclareMathOperator{\Im}{Im}
\DeclareMathOperator{\rk}{rk}
\DeclareMathOperator{\Tr}{Tr}
\DeclareMathOperator{\rad}{rad}
\DeclareMathOperator{\Ob}{Ob}
\DeclareMathOperator{\Mor}{Mor}
\DeclareMathOperator{\Sh}{Sh}
\newcommand{\G}{\mathsf G}
\newcommand{\BookChapter}[2]{\setcounter{chapter}{\numexpr#1-1\relax}\chapter{#2}\setcounter{equation}{0}}
\newcommand{\Section}[2]{\section*{#1\quad #2}\addcontentsline{toc}{section}{#1\quad #2}}
\newcommand{\PartTitle}[2]{\bigskip\begin{center}{\color{Gold}\rule{18mm}{0.5pt}}\quad{\small\sffamily\color{Muted}#1}\quad{\color{Gold}\rule{18mm}{0.5pt}}\par\medskip{\large\sffamily\bfseries\color{Teal}#2}\end{center}\medskip}
\newcommand{\sourcepage}[2]{\par\hypertarget{source-#2}{}\label{source-#2}\marginpar{\raggedright\color{Muted}\sffamily\fontsize{7}{9}\selectfont 原书\par #1}\ignorespaces}
\DeclareRobustCommand{\TermUse}[3]{\texorpdfstring{\hypertarget{use-#2}{}\label{use-#2}\hyperlink{term-#1}{#3}}{#3}}
\pdfstringdefDisableCommands{\def\TermUse#1#2#3{#3}\def\mathfrak#1{#1}\def\mathbb#1{#1}\def\mathcal#1{#1}\def\G{G}\def\Spec{Spec}\def\Aff{Aff}\def\Sets{Sets}\def\Hom{Hom}}
\newcommand{\GlossHeading}[4]{\par\Needspace{5\baselineskip}\hypertarget{term-#1}{}\noindent{\large\bfseries\color{Ink}#2}\par\nopagebreak\noindent{\small\textsf{EN}\enspace #3\quad\textsf{FR}\enspace #4}\par\nopagebreak}
\newcommand{\GlossRefs}[1]{\noindent{\fontsize{8}{11}\selectfont\color{Muted}返回正文：#1}\par\medskip}
\newcommand{\TranslatorNote}[1]{\footnote{\textbf{译注：}#1}}
\begin{document}
\begin{titlepage}
\thispagestyle{empty}
\begin{tikzpicture}[remember picture,overlay]
\fill[Ink] (current page.south west) rectangle (current page.north east);
\begin{scope}[shift={(current page.south west)},x=1mm,y=1mm]
\clip (0,0) rectangle (180,255);
\draw[white,opacity=.055,step=8mm] (15,32) grid (188,173);
\foreach \t in {-4,-3,...,4}{
  \draw[Teal!45!white,opacity=.35,line width=.5pt,domain=-5:5,samples=180,smooth,variable=\x]
  plot ({100+16*\x},{104+1.25*(\x*\x*\x-5*\x)+\t*5});
}
\draw[Gold,line width=1.15pt,domain=-4.5:4.5,samples=170,smooth,variable=\x]
plot ({100+16*\x},{104+1.25*(\x*\x*\x-5*\x)});
\foreach \x/\y in {64/117,100/104,136/91}{\fill[Paper] (\x,\y) circle (1.0mm);\draw[Gold,opacity=.6] (\x,\y) circle (2.2mm);}
\draw[Gold,line width=.7pt] (23,215)--(42,215);
\node[anchor=west,text=Gold,font=\sffamily\fontsize{10}{12}\selectfont] at (46,215) {GROTHENDIECK\enspace /\enspace GAETA};
\node[anchor=north west,text=Paper,align=left,font=\sffamily\bfseries\fontsize{30}{39}\selectfont] at (22,202) {函子代数几何\\导论};
\node[anchor=west,text=Paper,font=\sffamily\fontsize{16}{20}\selectfont] at (23,163) {第一卷\quad 仿射代数几何};
\node[anchor=west,text=Paper!65!Ink,align=left,font=\fontsize{9}{14}\selectfont] at (23,147) {INTRODUCTION TO FUNCTORIAL\\ALGEBRAIC GEOMETRY\\[2pt]VOL. I\enspace AFFINE ALGEBRAIC GEOMETRY};
\node[text=Paper!85!Ink,font=\fontsize{14}{17}\selectfont] (a) at (42,57) {$\mathfrak X_A$};
\node[text=Paper!85!Ink,font=\fontsize{14}{17}\selectfont] (b) at (137,57) {$\Hom_{k\text{-}\mathrm{alg}}(A,-)$};
\draw[-{Stealth[length=2mm]},Paper!60!Ink,line width=.5pt] (a)--node[above,text=Gold,font=\small] {$\sim$} (b);
\draw[Gold,opacity=.6,line width=.4pt] (23,37)--(157,37);
\node[anchor=west,text=Paper!80!Ink,font=\sffamily\fontsize{9}{12}\selectfont] at (23,27) {A. Grothendieck 讲授\qquad F. Gaeta 笔记};
\node[anchor=west,text=Paper!55!Ink,font=\sffamily\fontsize{8}{10}\selectfont] at (23,19) {中文译本\quad /\quad 中·英·法术语互链};
\end{scope}
\end{tikzpicture}
\null
\end{titlepage}

\frontmatter
\thispagestyle{plain}
\vspace*{18mm}
\noindent{\sffamily\color{Teal}\large 中文译本说明}\par\bigskip
本书据 A. Grothendieck 的暑期讲课、由 Federico Gaeta 整理的 \textit{Introduction to Functorial Algebraic Geometry, Vol. I: Affine Algebraic Geometry} 扫描本译出。原书由纽约州立大学布法罗分校数学系印行。译文保留原书的章节、命题与公式编号，并译出原书中的补充说明和脚注。

正文外侧的灰色页码对应原书印刷页码，供读者回查扫描本。原书中的明显排印问题在译注中说明。数学记号按原书重排，交换图以矢量形式绘制。

正文中的青色技术术语可点击跳至书末的中英法对照表。每个术语条目下均列有返回正文的链接；数字为本译本页码，同页出现多次时，括号内的序号分别链接到该页的各个出现位置。PDF 阅读器的“返回上一视图”功能也可直接返回刚才的阅读位置。

术语表中的法语采用代数几何及范畴论的惯用术语。原书的英文表达与现代常用说法有别时，以正文所指的数学概念为准。

\vfill
\noindent{\color{Gold}\rule{\linewidth}{0.5pt}}\par\smallskip
\noindent{\small\color{Muted}讲授：A. Grothendieck\qquad 笔记：F. Gaeta\par 排版：XeLaTeX\quad\textbullet\quad 中文译本}
\clearpage
\tableofcontents
\mainmatter
\input{build/linked-a}
\input{build/linked-b}
\input{build/linked-c}
\input{build/linked-d}
\input{build/linked-e}
\input{build/linked-f}
\input{build/linked-g}
\backmatter
\chapter{中英法技术名词对照表}
\label{glossary}
\noindent 本表按中文术语的阅读主题编排。点击下方页码或同页序号，可返回正文中相应的出现位置。\par\bigskip
\input{build/glossary}
\end{document}
